%% ============================================================================
\section{Βελτιστοποίηση Decode (Reed-Muller + Reed-Solomon)}
\label{sec:decode-optimization}

Το υποσύστημα αποκωδικοποίησης εκτελείται αποκλειστικά κατά το \Decaps και
αποτελείται από δύο διαδοχικά στάδια: τον εσωτερικό κώδικα \gls{rm}(1,7)
(bit-level διόρθωση εντός κάθε symbol) και τον εξωτερικό κώδικα \gls{rs}
πάνω από $\GF{2^8}$ (symbol-level διόρθωση μεταξύ symbols).  Μετά τη
βελτιστοποίηση του πολυωνυμικού πολλαπλασιασμού (\cref{ch:polymul-optimization})
και του Keccak (\cref{sec:keccak-optimization}), το decode pipeline αποτελεί
τον τρίτο μεγαλύτερο πυρήνα — καταναλώνοντας \pct{21}–\pct{34} του
\Decaps \instret στο scalar reference ανάλογα με το security level.

Η βελτιστοποίηση ακολουθεί τρία tiers αυξανόμενης ISA
πολυπλοκότητας: scalar με \Zbc injection (Tier~1), Base-\RVV
διανυσματοποίηση \gls{rm} και \gls{rs} syndromes (Tier~2), και πλήρης
\Zvbc vectorization όλων των αξόνων \gls{rs} (Tier~3).  Η τελική Tier~3
υλοποίηση επιτυγχάνει \speedup{6.4}–\speedup{7.3} στο decode pipeline
σε σχέση με το scalar reference.

%% --------------------------------------------------------------------------
\subsection{Scalar Path: \Zbc Injection στη \code{gf\_mul}}
\label{sec:decode-scalar}

Η μοναδική βελτιστοποίηση στο scalar path αφορά τον πυρήνα αριθμητικής
$\GF{2^8}$.  Η reference υλοποίηση (\code{gf\_carryless\_mul}) χρησιμοποιεί
έναν constant-time schoolbook αλγόριθμο που απαιτεί $\approx\num{15}$
εντολές ανά $8 \times 8$-bit carry-less γινόμενο.  Στο Tier~1, η
αντικατάσταση με μία μόνο εντολή \clmul (μέσω \Zbc intrinsic) μειώνει
το κόστος σε μία εντολή για το γινόμενο, ακολουθούμενη από $\approx\num{12}$
εντολές modular reduction μέσω \code{gf\_reduce} — σύνολο
$\approx\num{14}$ εντολές ανά κλήση \code{gf\_mul}, έναντι
$\approx\num{27}$ στο software path (\cref{lst:gf-mul-zbc}).

\begin{lstlisting}[language=C, caption={Scalar \code{gf\_mul} με \Zbc injection.},
  label={lst:gf-mul-zbc}]
uint16_t gf_mul(uint16_t a, uint16_t b) {
    uint64_t prod = __riscv_clmul_64(
        (uint64_t)(a & 0xFF), (uint64_t)(b & 0xFF));
    return gf_reduce((uint16_t)prod);
}
\end{lstlisting}

Η βελτίωση εφαρμόζεται οριζόντια σε ολόκληρο τον αποκωδικοποιητή: κάθε
κλήση \code{gf\_mul} στους υπολογισμούς syndromes, \gls{bm}, Forney, και
\gls{fft} root-finding ωφελείται ομοιόμορφα.  Η μείωση ανά κλήση
($\approx\pct{48}$) αθροίζεται σε χιλιάδες κλήσεις ανά \Decaps — ιδιαίτερα
στο HQC-256 όπου $\delta = 29$ και οι εσωτερικοί βρόχοι εκτελούν
$\BigO{\delta^2}$ πολλαπλασιασμούς.

Η \Zbb δεν συνεισφέρει στο decode path: ο \gls{rm} hadamard χρησιμοποιεί
πρόσθεση/αφαίρεση \code{int16}, ενώ η \code{gf\_reduce} χρησιμοποιεί
ολισθήσεις και XOR — κανένα pattern δεν αντιστοιχεί σε \Zbb εντολές
(rotation, \insn{andn}, count-leading-zeros).

%% --------------------------------------------------------------------------
\subsection{Base-RVV Path: Διανυσματοποίηση RM και RS Syndromes}
\label{sec:decode-rvv}

Το Tier~2 εισάγει πλήρη \RVV vectorization του \gls{rm} αποκωδικοποιητή
και vectorized syndrome computation για τον \gls{rs}.  Οι υπόλοιποι άξονες
\gls{rs} (B, C, D) παραμένουν scalar.  Η ISA απαίτηση περιορίζεται στη
βασική V extension — δεν απαιτείται \Zvbc.

\paragraph{Reed-Muller: Vectorized FHT μέσω Strided Loads.}
Ο αποκωδικοποιητής RM(1,7) εκτελεί τρεις υπο-συναρτήσεις ανά symbol:
\code{expand\_and\_sum} (εξαγωγή bits από MULTIPLICITY αντίγραφα),
\code{hadamard} (7-pass Fast Hadamard Transform σε \num{128} στοιχεία),
και \code{find\_peaks} (εύρεση θέσης μεγίστου).

Ο πυρήνας της vectorization είναι ο FHT.  Κάθε pass εκτελεί \num{64}
butterfly ζεύγη με stride-2 πρόσβαση — pattern που αντιστοιχεί απευθείας
στην εντολή \insn{vlse16} (strided load).  Στο \VLEN=256, SEW=16, LMUL=4,
η hardware vector length φτάνει \num{64} στοιχεία — επαρκή για κάλυψη
όλων των butterfly ζευγών σε μία μόνο vector πράξη (\cref{lst:hadamard-rvv}):

\begin{lstlisting}[language=C, caption={Vectorized FHT butterfly pass.},
  label={lst:hadamard-rvv}]
size_t vl = __riscv_vsetvl_e16m4(64);
vint16m4_t even = vlse16_v(in,     stride=4); // p1[0,2,...,126]
vint16m4_t odd  = vlse16_v(in + 1, stride=4); // p1[1,3,...,127]
vse16_v(out,      vadd_vv(even, odd, vl));     // sums
vse16_v(out + 64, vsub_vv(even, odd, vl));     // diffs
\end{lstlisting}

Κάθε pass απαιτεί ακριβώς \num{6} vector εντολές (2 strided loads, 2
αριθμητικές, 2 contiguous stores) — σύνολο \num{42} vector εντολές για
τους \num{7} passes.  Η scalar reference εκτελεί \num{896} αριθμητικές
πράξεις (πρόσθεση/αφαίρεση/αποθήκευση) για τον ίδιο υπολογισμό.

Η \code{expand\_and\_sum} αξιοποιεί την εντολή \insn{vid} για δημιουργία
shift amounts $[0, 1, \ldots, 31]$ σε SEW=32, LMUL=4 ($\text{vl}=32$),
εξάγοντας όλα τα \num{32} bits ενός \code{uint32} word σε μία vector
πράξη μέσω \insn{vsrl.vv} + \insn{vand}.  Η \code{find\_peaks} χρησιμοποιεί
LMUL=8 ($\text{vl}=128$) για σάρωση ολόκληρου του μετασχηματισμένου
πίνακα σε ένα pass: \insn{vredmaxu} για εύρεση μεγίστου, \insn{vmseq} +
\insn{vfirst\_m} για εντοπισμό θέσης.

\paragraph{Reed-Solomon Axis A: Vectorized Syndrome Computation.}
Κάθε ένα από τα $2\delta$ syndromes αποτελεί ανεξάρτητο εσωτερικό γινόμενο
μήκους $N_1 - 1$ πάνω από $\GF{2^8}$.  Ο εξωτερικός βρόχος (ανά syndrome)
παραμένει scalar, ενώ ο εσωτερικός βρόχος (ανά στοιχείο) vectorized πλήρως.

Η στρατηγική αξιοποιεί την ταυτότητα λογαρίθμου/αντιλογαρίθμου:
\begin{equation}
  \text{gf\_mul}(a, b) = \text{gf\_exp}\bigl[(\text{gf\_log}[a] + \text{gf\_log}[b]) \bmod 255\bigr]
\end{equation}
%
μέσω πέντε vectorized βημάτων ανά iteration: (1)~indexed gather στον
\code{gf\_log[]} πίνακα μέσω \insn{vluxei16}, (2)~modular πρόσθεση
λογαρίθμων με branchless conditional subtract ($\geq 255 \Rightarrow -255$),
(3)~indexed gather στον \code{gf\_exp[]} πίνακα, (4)~zero-masking μέσω
\insn{vmsne} + \insn{vmand} + \insn{vmerge} για τη συνθήκη $0 \cdot x = 0$,
και (5)~XOR-reduction μέσω \insn{vredxor}.

Η παραμετροποίηση SEW=16, LMUL=4 παρέχει $\text{vl}=64$ lanes στο
\VLEN=256 — \num{4} φορές ευρύτερο παραλληλισμό ανά pass σε σχέση με
την \Zvbc lane-packing προσέγγιση (SEW=64, $\text{vl}=16$).  Ο αριθμός
passes ανά syndrome είναι $\lceil (N_1 - 1) / 64 \rceil$: μόλις \num{1}
pass για HQC-128/192 ($N_1 = 46, 56$) και \num{2} passes για HQC-256
($N_1 = 90$).  Η υλοποίηση είναι \gls{vla}: χρησιμοποιεί
\code{vsetvl\_e16m4} για ερώτηση του hardware vector length, εξασφαλίζοντας
ορθότητα σε οποιοδήποτε $\VLEN \geq 64$.

Το trade-off έναντι της \Zvbc προσέγγισης (Track~C) συνοψίζεται ως εξής:
η RVV gather μέθοδος εμφανίζει υψηλότερο εντολικό κόστος ανά στοιχείο
($\approx\num{12}$ εντολές) αλλά επεξεργάζεται \num{4} φορές περισσότερα
στοιχεία ανά pass.  Επιπλέον, δεν απαιτεί Vector-Crypto extensions —
μόνο τη βασική V extension και scalar \Zbc για τα fallback axes B/C/D.
Σε υλοποιήσεις με αποδοτικές indexed-load μονάδες, ο ευρύτερος
παραλληλισμός υπερτερεί.

%% --------------------------------------------------------------------------
\subsection{Vector-Crypto Path: \Zvbc RS (Berlekamp-Massey, Forney)}
\label{sec:decode-zvbc-bm-forney}

Το Tier~3 επεκτείνει το Tier~2 vectorizing τους εναπομείναντες εσωτερικούς
βρόχους \gls{rs} (Axes B και D) μέσω \Zvbc \insn{vclmul.vv} lane-packing.
Η στρατηγική τοποθετεί ένα στοιχείο $\GF{2^8}$ ανά 64-bit lane, εκτελεί
carry-less πολλαπλασιασμό μέσω \insn{vclmul.vv}, και εφαρμόζει inline
Barrett reduction — εξαλείφοντας πλήρως τις scalar \code{gf\_mul} κλήσεις
στους εσωτερικούς βρόχους.

\paragraph{Axis B: Berlekamp-Massey (\code{compute\_elp\_zvbc}).}
Ο αλγόριθμος \gls{bm} εκτελεί $2\delta$ εξωτερικές επαναλήψεις, κάθε μία
εξαρτώμενη διαδοχικά από την τιμή discrepancy $d$ και τις μεταβλητές
ελέγχου βαθμού.  Ο εξωτερικός βρόχος δεν παραλληλοποιείται.  Ωστόσο, κάθε
επανάληψη περιέχει τρεις εσωτερικούς βρόχους αυξανόμενου μήκους (από
\num{1} έως $\delta$), οι οποίοι είναι ανεξάρτητοι εντός κάθε επανάληψης:

\begin{enumerate}
  \item \textbf{Sigma SAXPY}: $\sigma[i] \mathrel{\oplus}= dd \cdot X_{\sigma_p}[i]$
        — broadcast του scalar $dd$ μέσω \insn{vmv.v.x}, \insn{vclmul.vv}
        με τα zero-extended $X_{\sigma_p}$ στοιχεία, Barrett reduction, και
        XOR με τα τρέχοντα $\sigma$ values.
  \item \textbf{Blend}: constant-time conditional copy μέσω bitwise select
        pattern (\insn{vand} + \insn{vxor}) σε SEW=16, LMUL=4 — χωρίς
        $\GF{2^8}$ πολλαπλασιασμό.
  \item \textbf{Discrepancy dot-product}: φόρτωση $\sigma[1..n]$ contiguously
        και $S[\mu..\mu+1-n]$ σε αντίστροφη σειρά μέσω \insn{vlse16} με
        αρνητικό stride $-2$ bytes, \insn{vclmul.vv}, reduction, και
        XOR-accumulation μέσω \insn{vredxor}.
\end{enumerate}

Στο LMUL=4, SEW=64, $\text{vl}=16$ (\VLEN=256): οι εσωτερικοί βρόχοι
χωράνε σε \num{1} vector pass για HQC-128 ($\delta=15$) και HQC-192
($\delta=16$), και σε \num{2} passes για HQC-256 ($\delta=29$).

\paragraph{Axis D: Forney Error-Values (\code{compute\_error\_values\_zvbc}).}
Ο τύπος Forney υπολογίζει τιμές σφάλματος $e_j$ για κάθε θέση σφάλματος.
Δύο δομές vectorized:

\begin{itemize}
  \item \textbf{Αριθμητής (dot-product)}: μετά τη scalar προ-υπολογισμό
        δυνάμεων \code{inv\_pow[1..$\delta$]}, το εσωτερικό γινόμενο με
        $z[1..\delta]$ vectorized μέσω \insn{vclmul.vv} + \insn{vredxor} —
        δομικά ταυτόσημο με τον Axis~A syndrome dot-product.
  \item \textbf{Παρονομαστής (inner-k batch)}: ο modular δείκτης
        $\beta_j[(i+k) \bmod \delta]$ αποσυντίθεται σε δύο συνεχόμενα
        τμήματα μνήμης (αποφεύγοντας indexed gather), κάθε ένα
        επεξεργαζόμενο μέσω broadcast \insn{vmv.v.x} + \insn{vclmul.vv} +
        Barrett reduction + \insn{vxor.vx} (XOR με σταθερά \num{1}).  Το
        τελικό sequential γινόμενο $\text{tmp2} = \prod_{k} \text{inner\_k}[k]$
        παραμένει scalar — η $\BigO{\delta}$ αλυσίδα δεν παραλληλοποιείται
        χωρίς prefix-product scan, ο οποίος δεν αποσβένεται στα μικρά
        $\delta$ του \HQC.
\end{itemize}

Ο \cref{tab:decode-rs-zvbc} παρουσιάζει τα αποτελέσματα της vectorization
Axes B+D σε απομόνωση (πριν την ενσωμάτωση του Axis~C).

\begin{table}[htbp]
  \centering
  \caption{Αποτελέσματα \Zvbc vectorization Axes B+D στο \Decaps \instret.
           Η βελτίωση κλιμακώνεται με $\delta^2$ — μέγιστη στο HQC-256
           ($\delta=29$).}
  \label{tab:decode-rs-zvbc}
  % RS Zvbc Axes B+D optimization results (decaps only)
\begin{tabular}{l S[table-format=7.0] S[table-format=7.0] S[table-format=7.0] S[table-format=2.1]}
  \toprule
  {Security Level} & {Baseline \Decaps} & {Optimized \Decaps} & {$\Delta$ instret} & {$\Delta$ (\%)} \\
  \midrule
  HQC-128 & 521613  & 485507  & -36106  & -6.9 \\
  HQC-192 & 1150818 & 1110836 & -39982  & -3.5 \\
  HQC-256 & 1861032 & 1726034 & -134998 & -7.3 \\
  \bottomrule
\end{tabular}

\end{table}

Η κλιμάκωση επιβεβαιώνει τη θεωρητική πρόβλεψη: το HQC-256 ($\delta=29$)
εμφανίζει τη μεγαλύτερη απόλυτη μείωση ($-\cycles{134998}$ \instret,
$-\pct{7.3}$) διότι οι εσωτερικοί βρόχοι \gls{bm} και Forney εκτελούν
$\BigO{\delta^2}$ πολλαπλασιασμούς — πλήρως αξιοποιώντας τον vector
throughput.  Τα \KeyGen και \Encaps παραμένουν αμετάβλητα, επιβεβαιώνοντας
την απομόνωση της βελτιστοποίησης στο \Decaps path.

%% --------------------------------------------------------------------------
\subsection{Vector-Crypto Path: \Zvbc FFT Root-Finding (Axis C)}
\label{sec:decode-zvbc-fft}

Ο Axis~C αξιολογεί το error-locator πολυώνυμο $\sigma(x)$ σε όλα τα
$2^M = 256$ σημεία του $\GF{2^8}$ μέσω αναδρομικής radix αποσύνθεσης
(additive Gao-Mateer \gls{fft}).  Η reference υλοποίηση (\code{fft\_rec})
εκτελεί depth-first αναδρομική διάσχιση ενός δυαδικού δέντρου βάθους
\code{PARAM\_M} (= 8), με $\BigO{2^M}$ κλήσεις συναρτήσεων.

\paragraph{Μετασχηματισμός: Αναδρομικό → Επαναληπτικό BFS.}
Η κρίσιμη αρχιτεκτονική απόφαση ήταν η μετατροπή της αναδρομικής δομής
σε επαναληπτική breadth-first διάσχιση (\cref{fig:decode-fft-transform}).
Ο μετασχηματισμός αποτελείται από δύο passes:

\begin{itemize}
  \item \textbf{Pass 1 (top-down)}: για κάθε επίπεδο δέντρου από ρίζα
        προς φύλλα, εφαρμογή vectorized beta-scaling σε κάθε segment,
        ακολουθούμενη από radix αποσύνθεση σε δύο παιδιά.
  \item \textbf{Leaf base case}: vectorized XOR-broadcast για παραγωγή
        των $2^{m_{\text{leaf}}}$ εξόδων ανά φύλλο.
  \item \textbf{Pass 2 (bottom-up)}: butterfly combination παιδιών μέσω
        vectorized gammas-sums πολλαπλασιασμό.
\end{itemize}

Ο μετασχηματισμός εξαλείφει $\BigO{2^M}$ κλήσεις συναρτήσεων (function
call overhead: prologue/epilogue, stack frame allocation, argument passing)
και αντικαθιστά τις με explicit level-loops πάνω σε compile-time-bounded
arrays.  Η μνήμη εργασίας εκχωρείται στη στοίβα με σταθερά μεγέθη
καθοριζόμενα από \code{PARAM\_M} και \code{PARAM\_FFT} — χωρίς heap
allocation.

\begin{figure}[htbp]
  \centering
  % VARIANT (baseline refined, SAME layout) — not \input by main.tex; promote manually.
% decode-fft-transformation_v1: token migration + cleanup of the original
% side-by-side layout (raw red!8/green!8/blue!8 -> figthreat/figsolution/figprocess
% tokens; raw arrows -> figedge/*; \footnotesize -> figT/* scale). Central leaves
% (f01/f10) de-overlapped. The recursive tree is anchored so its root f(x) sits at
% the SAME height as the BFS "Level 0" box (root placed via (... |- lev0)), so the
% "Μετασχηματισμός" arrow is perfectly horizontal. Tree shape unchanged (calc
% offsets reproduce the liked layout). NOTE: keeps the wide side-by-side layout;
% use _vertical to fit \textwidth.
\tikzsetnextfilename{ch7-kernels-decode-fft-transformation_v1}
\begin{tikzpicture}[
  fftbox/.style ={figbox, align=center, minimum height=0.6cm},
  fftrec/.style ={fftbox, draw=figthreat,   fill=figthreatfill,   minimum width=1.6cm},
  fftit/.style  ={fftbox, draw=figsolution, fill=figsolutionfill, minimum width=4.6cm},
  fftcomb/.style={fftbox, draw=figprocess,  fill=figprocessfill,  minimum width=4.6cm},
]

% ===== Δεξιά: Επαναληπτικό (BFS) — defined first so the tree can align to lev0 =====
\node[figT/title] (rtitle) at (10.6,3.25) {Επαναληπτικό (BFS)};
\node[fftit,  below=4mm of rtitle] (lev0) {Level 0: 1 segment};
\node[fftit,  below=8mm of lev0]   (lev1) {Level 1: 2 segments};
\node[fftit,  below=8mm of lev1]   (lev2) {Level 2: 4 segments};
\node[fftit,  below=8mm of lev2]   (leaf) {Leaves: XOR-broadcast};
\node[fftcomb,below=8mm of leaf]   (comb) {Bottom-up: combine};

\draw[figedge] (lev0) -- (lev1) node[midway, right=3pt, figT/annot] {\insn{beta\_scale} + radix};
\draw[figedge] (lev1) -- (lev2) node[midway, right=3pt, figT/annot] {\insn{beta\_scale} + radix};
\draw[figedge] (lev2) -- (leaf) node[midway, right=3pt, figT/annot] {\insn{xor\_broadcast}};
\draw[figedge/weak] (leaf) -- (comb) node[midway, right=3pt, figT/annot] {\insn{gammas\_combine}};

\node[figT/callout, below=3mm of comb, align=center] {Σταθερός αριθμός επαναλήψεων βρόχου};

% ===== Αριστερά: Αναδρομικό (depth-first) — root aligned to lev0 (straight arrow) =====
% Root sits at lev0's height; rest of the tree is placed relative to root (same shape).
\node[fftrec] (root) at (2.775,0 |- lev0) {$f(x)$};
\node[fftrec] (l1)  at ($(root)+(-1.85,-1.20)$) {$f_0$};
\node[fftrec] (r1)  at ($(root)+( 1.85,-1.20)$) {$f_1$};
\node[fftrec] (l2a) at ($(root)+(-2.775,-2.40)$) {$f_{00}$};
\node[fftrec] (l2b) at ($(root)+(-0.925,-2.40)$) {$f_{01}$};
\node[fftrec] (r2a) at ($(root)+( 0.925,-2.40)$) {$f_{10}$};
\node[fftrec] (r2b) at ($(root)+( 2.775,-2.40)$) {$f_{11}$};
\node[figT/title] (ltitle) at ($(root)+(0,0.85)$) {Αναδρομικό (depth-first)};

\draw[figedge] (root) -- (l1) node[midway, left, figT/annot] {call};
\draw[figedge] (root) -- (r1) node[midway, right, figT/annot] {call};
\draw[figedge] (l1) -- (l2a);
\draw[figedge] (l1) -- (l2b);
\draw[figedge] (r1) -- (r2a);
\draw[figedge] (r1) -- (r2b);

\node[figT/callout, anchor=north, align=center] at ($(root)+(0,-3.05)$)
  {$\BigO{2^M}$ κλήσεις συναρτήσεων};

% ===== Μετασχηματισμός (οριζόντιο βέλος — root.east και lev0.west στο ίδιο y) =====
\draw[figedge/accent] ([xshift=0.5cm]root.east) -- ([xshift=-0.5cm]lev0.west)
  node[midway, above, figT/label] {Μετασχηματισμός};
\end{tikzpicture}

  \caption{Μετασχηματισμός αναδρομικού \code{fft\_rec} σε επαναληπτικό BFS.
           Αριστερά: depth-first αναδρομή με $\BigO{2^M}$ κλήσεις
           συναρτήσεων.  Δεξιά: level-by-level βρόχος με σταθερό αριθμό
           επαναλήψεων και vectorized helpers σε κάθε επίπεδο.}
  \label{fig:decode-fft-transform}
\end{figure}

\paragraph{Τρεις Vectorized Helpers.}
Η επαναληπτική δομή ενεργοποιεί τρεις vectorized πυρήνες:

\begin{enumerate}
  \item \textbf{\code{beta\_scale\_zvbc}$(f, n, \beta)$}: υπολογίζει
        $f[i] = \text{gf\_mul}(\beta^i, f[i])$ για $i = 1, \ldots, n-1$.
        Οι δυνάμεις $\beta^i$ προ-υπολογίζονται σε bounded stack array.
        Η vectorized πράξη χρησιμοποιεί το pipeline: \insn{vle16} →
        \insn{vzext} u16→u64 → \insn{vclmul.vv} → Barrett reduction →
        narrowing u64→u16 → \insn{vse16}.  Η περίπτωση $\beta = 1$
        (identity) εκτελείται χωρίς branch — η \code{gf\_mul}(1, x) = x
        αποτελεί fixed-latency vector no-op.
  \item \textbf{\code{xor\_broadcast\_zvbc}$(dst, src, len, val)$}: υπολογίζει
        $dst[k] = src[k] \oplus val$ μέσω \gls{vla} loop:
        \insn{vle16}/\insn{vxor.vx}/\insn{vse16} σε SEW=16, LMUL=4
        ($\text{vl}=64$ στο \VLEN=256).  Αντικαθιστά τον scalar βρόχο
        XOR-broadcast σε κάθε επίπεδο δέντρου.
  \item \textbf{\code{gammas\_sums\_combine\_zvbc}$(out, u, gs, v, count)$}:
        υπολογίζει $out[i] = u[i] \oplus \text{gf\_mul}(gs[i], v[i])$ μέσω
        του πλήρους \insn{vclmul.vv} + Barrett reduce + \insn{vxor.vv}
        pipeline.  Μία scalar-broadcast παραλλαγή χειρίζεται την περίπτωση
        $\text{sc} \leq 3$ όπου το $v$ είναι σταθερά.
\end{enumerate}

\paragraph{Αποτελέσματα.}
Ο \cref{tab:decode-instret} παρουσιάζει τη συνολική επιτάχυνση του decode
pipeline (πλήρες RM + RS, συμπεριλαμβανομένου του Axis~C) σε σχέση με
το scalar reference.  Ο \cref{tab:decode-axisc} απομονώνει τη συνεισφορά
του Axis~C μέσω διαφοράς με το προηγούμενο Zvbc baseline (ίδιοι Axes
A/B/D, scalar \code{fft\_rec} για Axis~C).

\begin{table}[htbp]
  \centering
  \caption{Συνολικό decode pipeline \instret: scalar reference έναντι
           πλήρους \Zvbc υλοποίησης (Axes A+B+C+D).  Μέτρηση σε Spike,
           \VLEN=256.}
  \label{tab:decode-instret}
  % Decode pipeline instret comparison table
\begin{tabular}{l S[table-format=7.0] S[table-format=6.0] S[table-format=1.2]}
  \toprule
  {Security Level} & {Scalar reference} & {Zvbc (πλήρες pipeline)} & {Speedup} \\
  \midrule
  HQC-128 & 384986  & 60353  & 6.38 \\
  HQC-192 & 562125  & 77114  & 7.29 \\
  HQC-256 & 1059605 & 159206 & 6.66 \\
  \bottomrule
\end{tabular}

\end{table}

\begin{table}[htbp]
  \centering
  \caption{Απομονωμένη συνεισφορά Axis~C (FFT root-finding): σύγκριση
           \Zvbc pipeline με scalar FFT έναντι vectorized FFT.}
  \label{tab:decode-axisc}
  % Axis C (FFT root-finding) isolated contribution table
\begin{tabular}{l S[table-format=6.0] S[table-format=6.0] S[table-format=2.1]}
  \toprule
  {Security Level} & {Zvbc (scalar FFT)} & {Zvbc (vectorized FFT)} & {$\Delta$ (\%)} \\
  \midrule
  HQC-128 & 88011  & 56326  & -36.0 \\
  HQC-192 & 119982 & 72012  & -40.0 \\
  HQC-256 & 204354 & 146513 & -28.3 \\
  \bottomrule
\end{tabular}

\end{table}

Η vectorized FFT επιτυγχάνει $-\pct{31.4}$ (HQC-128) έως $-\pct{35.7}$
(HQC-192) μείωση στο Axis~C component.  Η χαμηλότερη βελτίωση στο HQC-256
($-\pct{22.1}$) οφείλεται στο μεγαλύτερο \code{PARAM\_FFT} (= 5 αντί 4),
το οποίο αυξάνει τον αριθμό επιπέδων δέντρου και τα μεγέθη segments —
μειώνοντας αναλογικά το ποσοστό function-call overhead που εξαλείφθηκε.

Συνολικά, η πλήρης Tier~3 decode υλοποίηση (Axes A+B+C+D) επιτυγχάνει
\speedup{6.38} (HQC-128), \speedup{7.29} (HQC-192) και \speedup{6.66}
(HQC-256) στο decode pipeline.  Η ορθότητα επιβεβαιώθηκε μέσω
\num{300}/\num{300} \gls{kat} vectors σε όλα τα security levels.

%% --------------------------------------------------------------------------
\subsection{Βελτιστοποίηση GF($2^8$) Reduction: Single-Pass Barrett via \insn{vclmul}}
\label{sec:decode-barrett}

Η modular reduction $\bmod\; p(x) = x^8 + x^4 + x^3 + x^2 + 1$ (0x11D)
εφαρμόζεται μετά κάθε \insn{vclmul.vv} στο Zvbc decode pipeline.  Η αρχική
υλοποίηση χρησιμοποιούσε two-pass shift-and-XOR feedback (ταυτόσημη με τη
scalar \code{gf\_reduce}): \num{12} vector εντολές ανά reduction.  Η
βελτιστοποίηση αντικαθιστά αυτή τη δομή με single-pass Barrett reduction
χρησιμοποιώντας \insn{vclmul.vv} για εκτίμηση πηλίκου — μειώνοντας το
κόστος σε \num{6} vector εντολές.

\paragraph{Μαθηματική θεμελίωση.}
Η Barrett reduction στο $\GF{2}[x]$ υπολογίζει:
\begin{align}
  \text{high} &= v \gg 8 \label{eq:barrett-high} \\
  q &= (\text{high} \otimes \mu) \gg 8 \label{eq:barrett-q} \\
  \text{corr} &= q \otimes p(x) \label{eq:barrett-corr} \\
  r &= (v \oplus \text{corr}) \mathbin{\&} \text{0xFF} \label{eq:barrett-r}
\end{align}
%
όπου $\mu = \lfloor x^{16} / p(x) \rfloor = \text{0x11C}$ στο $\GF{2}[x]$
και $\otimes$ σημαίνει carry-less πολλαπλασιασμό.  Η κρίσιμη παρατήρηση
είναι ότι η χρήση $\text{0x11D}$ (αντί $\text{0x11C}$) ως $\mu$ παράγει
ταυτόσημα αποτελέσματα για όλες τις 14-bit εισόδους: η διαφορά $+1$ (bit~0)
δεν επηρεάζει το πηλίκο μετά τη δεξιά ολίσθηση κατά \num{8} θέσεις για
τελεστέους βαθμού $\leq 5$.  Η ιδιότητα αυτή (self-reciprocal) επιτρέπει
τη χρήση ενός μόνο constant vector και για τις δύο \insn{vclmul.vv} πράξεις.

\paragraph{Υλοποίηση.}
Η \code{gf\_reduce\_vec\_u64m4} (\cref{lst:barrett-single}) εκτελεί τη reduction σε \num{6} vector
εντολές (πλέον \num{1} \insn{vmv.v.x} broadcast που ανυψώνεται εκτός
βρόχου από τον μεταγλωττιστή):

\begin{lstlisting}[language=C, caption={Single-pass Barrett reduction μέσω \insn{vclmul.vv}.},
  label={lst:barrett-single}]
vuint64m4_t high = vsrl_vx(v, 8, vl);       // overflow [13:8]
vuint64m4_t q    = vclmul_vv(high, mu, vl);  // quotient est.
q                = vsrl_vx(q, 8, vl);        // shift quotient
vuint64m4_t corr = vclmul_vv(q, mu, vl);     // correction
v                = vxor_vv(v, corr, vl);      // apply
v                = vand_vx(v, 0xFF, vl);      // mask 8 bits
\end{lstlisting}

Η ορθότητα επαληθεύτηκε εξαντλητικά: όλες οι \num{16384} τιμές στο
εύρος $[0, \text{0x3FFF}]$ παράγουν ταυτόσημο αποτέλεσμα με τη reference
\code{gf\_reduce}.  Η βελτιστοποίηση εφαρμόζεται τόσο στο
\code{fft\_zvbc.c} όσο και στο \code{reed\_solomon\_zvbc.c}.

\paragraph{Εξερεύνηση \ext{Zvkg} (\insn{vgmul.vv}).}
Η εντολή \insn{vgmul.vv} (GHASH multiply) εξερευνήθηκε ως εναλλακτική
μονο-εντολής reduction.  Ωστόσο, η \ext{Zvkg} χρησιμοποιεί σταθερό
πολυώνυμο GHASH ($x^{128} + x^7 + x^2 + x + 1$) — ασύμβατο με το
$\GF{2^8}$ irreducible $\text{0x11D}$.  Πειραματική επαλήθευση επιβεβαίωσε
τη μη εφαρμοσιμότητα: \num{21}/\num{300} \gls{kat} failures (Option~A)
και \num{259}/\num{300} failures (Option~B).

\paragraph{Ποσοτικό αποτέλεσμα.}
Η single-pass Barrett μειώνει το \Decaps \instret κατά \cycles{2699}
($-\pct{0.67}$) στο HQC-128, χωρίς επίπτωση στα \KeyGen/\Encaps.  Η
μικρή απόλυτη τιμή αντανακλά το γεγονός ότι η reduction αποτελεί μικρό
κλάσμα του ήδη βελτιστοποιημένου decode pipeline — αλλά η \pct{50} μείωση
εντολών ανά reduction κλήση αθροίζεται σε χιλιάδες εκτελέσεις ανά \Decaps.
Η ορθότητα επιβεβαιώθηκε μέσω \num{300}/\num{300} \gls{kat} vectors.

